\section{Data and Evaluation}\label{sec:data-eval}

\subsection{Detectors and data}\label{sec:data}
Table~\ref{tab:data} lists the evaluation cells. A cell is one combination of
data, detector, and tracker (the trackers are introduced in
Sect.~\ref{sec:setup}). Three public datasets are used. MOT17
\cite{milan2016mot16} shows people in city scenes. As is common, each of its
training videos is split in time, and the second half, called val-half, is
used for evaluation. On MOT17 the trackers use the published detections of
the YOLOX-X detector \cite{ge2021yolox}, so every tracker sees exactly the
same boxes. A detector drops the boxes whose score is below a small value
before it passes them on; we call this value the emission floor. The two
published MOT17 detection sets have floors of 0.01 and 0.1
(Sect.~\ref{sec:tracker}). KITTI tracking \cite{geiger2012kitti} shows cars
and pedestrians seen from a driving car. On its training set the trackers run
with two detectors that differ from the MOT17 detector: YOLOv8n
\cite{jocher2023yolov8}, a small and fast model, and RT-DETR-L
\cite{zhao2024rtdetr}, a larger transformer model. UAVDT \cite{du2018uavdt}
shows vehicles seen from drones. Its test split, 20 sequences with 16,592
frames, was not used during the design; ByteTrack runs on it with the same two
detectors, restricted to their vehicle classes.

To study the score problem directly, the published MOT17 scores were also
changed before the tracker by four monotone transforms, which change the
score values but keep their order: cubing ($s^3$), halving ($0.5\,s$), and
logit temperatures $T=2$ and $T=0.5$, $\sigma(z/T)$, where $\sigma$ turns a
logit back into a score. A temperature divides the logit $z$ of each score by
$T$ before turning it back into a score: $T=2$ pulls all scores toward 0.5,
and $T=0.5$ pushes them toward 0 and 1. These transforms imitate the score
changes that a new or retrained detector can cause. The trackers kept their
published thresholds. The external systems use the detections released by
their authors on MOT17 val-half.

\begin{table*}[t]
\caption{Evaluation cells and their role. A cell is one combination of data,
detector, and tracker. Development cells were looked at while the layer was
designed; predeclared and held-out cells were chosen in advance and run once
with the frozen layer.}
\label{tab:data}
\centering\footnotesize
\setlength{\tabcolsep}{4pt}
\begin{tabular}{@{}>{\raggedright\arraybackslash}p{29mm}>{\raggedright\arraybackslash}p{25mm}>{\raggedright\arraybackslash}p{42mm}>{\raggedright\arraybackslash}p{25mm}@{}}
\toprule
Data & Detector & Trackers & Role\\
\midrule
MOT17 val-half (7 sequences) & published YOLOX-X detections & ByteTrack (2 settings), OC-SORT, BoostTrack, SparseTrack & development\\
KITTI tracking training (21 sequences) & RT-DETR-L; YOLOv8n & ByteTrack and BoT-SORT (RT-DETR-L); OC-SORT (YOLOv8n) & development (detector transfer)\\
MOT17 val-half, score transforms & published YOLOX-X, transformed scores & ByteTrack (2 settings), OC-SORT, BoostTrack & development (robustness)\\
MOT17 val-half & released PD-SORT / Hybrid-SORT detections & PD-SORT, Hybrid-SORT & predeclared external\\
UAVDT test (20 sequences) & RT-DETR-L; YOLOv8n & ByteTrack & predeclared held-out\\
KITTI tracking training, also with halved scores & RT-DETR-L; YOLOv8n & ByteTrack (RT-DETR-L); OC-SORT (YOLOv8n) & development (baselines, ACT, and components)\\
\bottomrule
\end{tabular}
\end{table*}


\subsection{Evaluation roles and metrics}\label{sec:roles}
Every cell has one of two roles (Table~\ref{tab:data}). Development cells
were looked at while \scl{} was designed: MOT17 with four trackers, KITTI
with two other detectors, and the score transforms. The other cells test
\scl{} on trackers and data that it had never seen. The two predeclared
external trackers, PD-SORT and Hybrid-SORT, were chosen and named before they
were run, and they were run once with the frozen configuration. The held-out
UAVDT test split was also run once with the frozen configuration, after a
written plan. To check that each external tracker was run correctly, our run
of the tracker alone is compared with the number reported by its authors and
marked EXACT (within 0.2 HOTA points) or CLOSE (within 1.0 HOTA point, with a
known cause that was not tuned).

\label{sec:metrics}Table~\ref{tab:metrics} gives the tracking metrics and
their definitions. They are built from four counts: false positives (FP),
boxes where there is no object; false negatives (FN), objects that were
missed; identity switches (IDS), cases where an object's ID changes; and the
number of ground-truth boxes (GT). HOTA is the main metric, because it
balances finding the objects and keeping their identities. MOTA counts errors
and reacts mostly to FP and FN, and the identity F1 score (IDF1) measures how
long each object keeps its correct ID. IDS should be read together with FN: a
tracker that follows more objects has more chances to switch identities, so
IDS can rise while the tracking as a whole becomes better. The metrics are
computed with the TrackEval library (commit 12c8791) and the official ground
truth of each benchmark, following MOTChallenge practice
\cite{dendorfer2021motchallenge}. On KITTI tracking, HOTA is the official
value averaged over cars and pedestrians. On UAVDT, all vehicles form one
class, and predictions inside the regions that the dataset marks as ignored
are removed.

\begin{table*}[t]
\caption{Tracking metrics used in this paper. GT is the number of
ground-truth boxes; FP, FN, and IDS are the false positives, missed objects
(false negatives), and identity switches, summed over all frames. For IDF1,
each ground-truth track is first matched to at most one output track, and
IDTP, IDFP, and IDFN are the true-positive, false-positive, and
false-negative boxes under this matching. For HOTA, at a localization
threshold $\alpha$, $\mathrm{TP}_\alpha$, $\mathrm{FN}_\alpha$, and
$\mathrm{FP}_\alpha$ are the matched, missed, and unmatched boxes, and
$\mathcal{A}(c)$ is the association accuracy of a matched pair $c$. HOTA,
MOTA, and IDF1 are given in percent.}
\label{tab:metrics}
\centering\footnotesize
\setlength{\tabcolsep}{4pt}
\renewcommand{\arraystretch}{1.6}
\begin{tabular}{@{}l>{\raggedright\arraybackslash}p{64mm}>{\raggedright\arraybackslash}p{46mm}@{}}
\toprule
Metric & Definition & What it measures\\
\midrule
FP, FN, IDS & counts summed over all frames & FP: boxes where there is no object; FN: objects that were missed; IDS: times an object's ID changes. Lower is better; IDS is read together with FN, because a tracker that follows more objects has more chances to switch IDs\\
MOTA \cite{bernardin2008clear} & $\displaystyle \mathrm{MOTA}=1-\frac{\mathrm{FN}+\mathrm{FP}+\mathrm{IDS}}{\mathrm{GT}}$ & one minus the share of errors; reacts mostly to false positives and missed objects\\
IDF1 \cite{ristani2016idf1} & $\displaystyle \mathrm{IDF1}=\frac{2\,\mathrm{IDTP}}{2\,\mathrm{IDTP}+\mathrm{IDFP}+\mathrm{IDFN}}$ & how long each object keeps its correct ID\\
HOTA \cite{luiten2021hota} & $\displaystyle \mathrm{HOTA}=\frac{1}{19}\sum_{\alpha}\sqrt{\mathrm{DetA}_\alpha\,\mathrm{AssA}_\alpha}$,\newline $\displaystyle \mathrm{DetA}_\alpha=\frac{|\mathrm{TP}_\alpha|}{|\mathrm{TP}_\alpha|+|\mathrm{FN}_\alpha|+|\mathrm{FP}_\alpha|}$,\quad $\displaystyle \mathrm{AssA}_\alpha=\frac{1}{|\mathrm{TP}_\alpha|}\sum_{c\in\mathrm{TP}_\alpha}\mathcal{A}(c)$,\newline with $\alpha=0.05,0.10,\dots,0.95$ & balance of finding the objects (detection accuracy, DetA) and keeping their identities (association accuracy, AssA); the main metric of this paper\\
\bottomrule
\end{tabular}
\end{table*}


All comparisons are paired: each tracker runs twice on the same sequences and
detections, once alone and once with \scl{}, and nothing else changes. The
difference is computed over all sequences together. To show how
certain a difference is, we use a bootstrap \cite{efron1993bootstrap}: the
sequences are drawn again at random, with repetition, 10,000 times (with the
fixed random seed 42, so that the result can be repeated),
the difference is computed again each time, and the middle 95\% of these
values gives the 95\% interval. A difference is called significant when its
interval does not include zero; bold intervals in the tables mark such
differences. As a simple check that needs no statistics, we also count the
sequences where HOTA rose, stayed the same ($|\Delta\mathrm{HOTA}|<0.01$), or
fell.

HOTA, MOTA, and IDF1 are given in percent, and higher values are better. FP,
FN, and IDS are counts, and lower values are better. Runtime is given in
milliseconds (ms) per frame, and lower values are better.
